\begin{Info}
% Universidade
{UNIVERSIDADE DO VALE DO ITAJAÍ}
% Escola
{ESCOLA POLITÉCNICA}
% Curso
{CURSO DE CIÊNCIA DA COMPUTAÇÃO}
% Titulo
{USO DE NERF PARA GERAÇÃO EM TEMPO REAL DE CENÁRIOS EM REALIDADE VIRTUAL: NERF, 3D GAUSSIAN SPLATTING E EXECUÇÃO ON-DEVICE NO META QUEST 3S}
% Autor
{Mateus Barbosa}
% Cidade e Data
{Itajaí (SC), Junho de 2026}
% Nome da Área de concentração
{Computação Gráfica}
% Orientador(a)
{Rodrigo Lyra}
% Coorientador(a) <Nome do Coorientador(a)>, <Titulação> %%%%%%%% Se não tiver coorientador deixe vazio

\end{Info}

%\begin{Dedicatoria}
% Dedicatória
% \end{Dedicatoria}

% \begin{Agradecimentos}
% Agradeço a todos.
% \end{Agradecimentos}

% \begin{Epigrafe}
% Epígrafe
% \end{Epigrafe}

\begin{Resumo}
BARBOSA, Mateus. Uso de NeRF para Geração em Tempo Real de Cenários em Realidade Virtual: NeRF, 3D Gaussian Splatting e Execução On-Device no Meta Quest 3S. Itajaí, 2026. Trabalho Técnico-científico de Conclusão de Curso (Graduação em Ciência da Computação) - Escola Politécnica, Universidade do Vale do Itajaí, Itajaí, 2026.

Este trabalho propõe um \textit{pipeline} de reconstrução de cenários 3D fotorrealistas para Realidade Virtual autônoma, articulando \textit{Neural Radiance Fields} (NeRF) e \textit{3D Gaussian Splatting} (3D-GS). A proposta parte do problema de que métodos neurais de reconstrução alcançam elevada fidelidade visual, mas ainda apresentam limitações práticas quando precisam operar em tempo real, com baixa latência e sob restrições de hardware móvel, como ocorre no Meta Quest 3S. A metodologia prevista envolve a captura de imagens RGB de cenas estáticas, o pré-processamento e a organização dos dados, o treinamento de um modelo NeRF como linha de base visual, a conversão ou reconstrução da cena em 3D-GS e a integração da representação resultante em um ambiente de Realidade Virtual com Unreal Engine. A avaliação proposta considera métricas de qualidade visual, como PSNR, SSIM e LPIPS, e métricas de desempenho, como FPS, tempo de quadro e uso de memória. Como contribuição, o trabalho estrutura uma base conceitual e metodológica para investigar o compromisso entre fidelidade visual e desempenho computacional na visualização imersiva de cenas reconstruídas em dispositivos de VR autônomos.

\textbf{Palavras-chave}: NeRF, 3D Gaussian Splatting, Realidade Virtual, Meta Quest 3S, Computação Gráfica.
\end{Resumo}

\begin{Abstract}
BARBOSA, Mateus. Use of NeRF for Real-Time Generation of Virtual Reality Scenes: NeRF, 3D Gaussian Splatting and On-Device Execution on Meta Quest 3S. Itajaí, 2026. Undergraduate Final Project (Bachelor's Degree in Computer Science) - Polytechnic School, Universidade do Vale do Itajaí, Itajaí, 2026.

This work proposes a pipeline for reconstructing photorealistic 3D scenes for standalone Virtual Reality, combining \textit{Neural Radiance Fields} (NeRF) and \textit{3D Gaussian Splatting} (3D-GS). The proposal is motivated by the fact that neural reconstruction methods can achieve high visual fidelity, but still face practical limitations when real-time execution, low latency, and mobile hardware constraints are required, as in the Meta Quest 3S. The planned methodology includes capturing RGB images of static scenes, preprocessing and organizing the data, training a NeRF model as a visual baseline, converting or reconstructing the scene into 3D-GS, and integrating the resulting representation into a Virtual Reality environment using Unreal Engine. The proposed evaluation considers visual quality metrics, such as PSNR, SSIM, and LPIPS, as well as performance metrics, including FPS, frame time, and memory usage. As its contribution, this work structures a conceptual and methodological foundation to investigate the trade-off between visual fidelity and computational performance in immersive visualization of reconstructed scenes on standalone VR devices.

\textbf{Keywords}: NeRF, 3D Gaussian Splatting, Virtual Reality, Meta Quest 3S, Computer Graphics.
\end{Abstract}
